\documentclass[11pt]{article}
\usepackage{amsmath}
\usepackage{amssymb}
\usepackage{amsthm}
\usepackage{enumitem}
\usepackage[margin=1in]{geometry}
\usepackage{multicol}
\usepackage{soul}
\usepackage{tikz}

% Load hyperref to automatically generate the PDF outline from structural tags
\usepackage[pdftex, bookmarks=true, bookmarksnumbered=true]{hyperref}

% Optional: Use the bookmark package for faster rendering and advanced control
\usepackage{bookmark} 

\newenvironment{case}[2]{\par\textbf{Case #1: #2.}}{}
\title{Homework \#4\\{\large Math 104: Intro to Analysis \\ Summer 2026} \\ {\normalsize UC Berkeley}}
\date{July 24, 2026}
\author{Donald Aingworth IV}

\begin{document}
\renewcommand\qedsymbol{OE$\Delta$}
\newcommand{\abs}[1]{\left| #1 \right|}
\newcommand{\andtxt}{\text{ and }}
\newcommand{\Ais}{\ensuremath{\overset{A}{=}}}
\newcommand{\Cis}{\ensuremath{\overset{C}{=}}}
\newcommand{\closure}[1]{\ensuremath{\overline{#1}}}
\newcommand{\encircle}[1]{%
  \tikz[baseline=(char.base)]{%
    \node[shape=circle, draw, inner sep=1pt] (char) {#1};%
  }%
}
\newcommand{\mathif}{\text{if }}
\newcommand{\seq}[2][n]{\ensuremath{\left\{ #2 _{#1} \right\}}}
\newcommand{\zz}[1]{$\mathbb{Z}/ #1 \mathbb{Z}$}

\maketitle

\tableofcontents

\pagebreak
\setcounter{section}{-1}
\section{Various Theorems}

\pagebreak
\section{Problem 1}
    Let $(X, d_X)$ and $(Y, d_Y)$ be metric spaces. A function $f : (X, d_X) \to (Y, d_Y)$ is called \textit{Lipschitz continuous} if there is a constant $L > 0$ so that
    \begin{equation*}
        d_Y(f(p), f(q)) \leq Ld_X(p, q)
    \end{equation*}
    for all $p, q \in X$. For $\alpha > 0$, we say $f$ is \textit{H\"older continuous (with exponent $\alpha$)} if there is a constant $C > 0$ so that
    \begin{equation*}
        d_Y(f(p), f(q)) \leq C d_X(p, q)^\alpha
    \end{equation*}
    for all $p, q \in X$. In particular, H\"older continuous with exponent $1$ means the same thing as Lipschitz.
    \begin{enumerate}[label=(\alph*)]
        \item 
        Fix a point $p \in X$. Show that the map $g : X \to \mathbb{R}$ defined by $g(x) = d(x, p)$ is Lipschitz continuous, with constant $L = 1$.
        
        \item 
        Prove that any H\"older continuous function (with any exponent $\alpha > 0$) is uniformly continuous.

        \item 
        Prove that if $f$ is Lipschitz continuous and bounded, then $f$ is H\"older continuous with exponent $\alpha$, for any $0 < \alpha \leq 1$.

        \item 
        Let $[a, b]$ be a closed interval, and suppose $f : [a, b] \to \mathbb{R}$ is $C^1$ (meaning its first derivative exists and is continuous on $[a, b]$). Show that $f$ is Lipschitz continuous.

        \item 
        Suppose $f : \mathbb{R} \to \mathbb{R}$ is H\"older continuous with exponent $\alpha > 1$. Show that $f$ is constant.
        
    \end{enumerate}

\pagebreak
\section{Problem 2}
    On the interval $(100, \infty)$, define functions $f(x) = x + \sin(x)$, $g(x) = x$, $p(x) = x + \sin(x)\cos(x)$, and $q(x) = p(x)e^{\sin(x)}$.
    \begin{enumerate}[label=(\alph*)]
        \item 
        Using facts we proved about limits and derivatives, prove that
        \begin{equation*}
            \lim_{x \to \infty} \frac{f(x)}{g(x)} = 1, \quad \lim_{x \to \infty} \frac{f'(x)}{g'(x)} \;\; \text{DNE}, \quad \lim_{x \to \infty} \frac{p(x)}{q(x)} \;\; \text{DNE}, \text{ and } \lim_{x \to \infty} \frac{p'(x)}{q'(x)} = 0.\\
        \end{equation*}
        \item 
        For both scenarios $\frac{f(x)}{g(x)}$ and $\frac{p(x)}{q(x)}$, explain why these results do not contradict L'H\^opital's rule.
        
    \end{enumerate}

    \textit{Note: even though we haven't learned about $\sin$ and $\cos$ yet, you are free to use any familiar properties of them you wish.}
\end{document}